In this section we will recover several known results, some in greater generality than in their original statements, dealing with different algebraic structures on classes of continuous functions and their isomorphisms.%. Moreover, the order provided from \ref{subsectionkaplansky}-\ref{subsectionbanachstone} allows us to apply each result in order to obtain the next one.
%
%In this section and the next, we will consider different algebraic signatures on $H_X$ and $H_Y$, and isomorphisms $T:\mathcal{A}(X)\to\mathcal{A}(Y)$ for appropriate $\mathcal{A}(X)$ and $\mathcal{A}(Y)$.

The general procedure we will use is the following: Suppose that $X$ and $Y$ are locally compact Hausdorff spaces, $H_X$ and $H_Y$ are Hausdorff spaces, $\theta_X\in C(X,H_X)$, $\theta_Y\in C(Y,H_Y)$ and $\mathcal{A}(X)$ and $\mathcal{A}(Y)$ regular subsets of $C_c(X,\theta_X)$ and $C_c(Y,\theta_Y)$, respectively.
\begin{enumerate}\label{generalprocedureforclassification}
\item Describe the relation $\perpp$ with the algebraic structure at hand. In general, we will first describe $\perp$ and use it in order to describe $\perpp$.
\item Given an algebraic isomorphism $T:\mathcal{A}(X)\to\mathcal{A}(Y)$ for appropriate classes of functions $\mathcal{A}(X)$ and $\mathcal{A}(Y)$, the first item ensures that $T$ is a $\perpp$-isomorphism, so let $\phi:Y\to X$ be the $T$-homeomorphism.
\item Prove that $T$ is $\phi$-basic. Let $\chi$ be the $(\phi,T)$-transform.
\item Items 1.-3.\ also apply to $T^{-1}$, which is thus also a basic $\perpp$-isomorphism.
\item Items 3.\ and 4.\  imply, by Proposition \ref{propositionbasicisomorphisms}(a) and (b), that the sections $\chi(y,\cdot)$ are injective and in the case that $\mathcal{A}(Y)$ is regular, item (d) implies that $\chi(y,\cdot)$ is also surjective.
\item If we have a classification of algebraic isomorphisms between $H_X$ and $H_Y$, this classification will apply to each section $\chi(y,\cdot)$ of the $(\phi,T)$-transform, by Proposition \ref{propositionmodelmorphism}. This will in turn describe $T$ completely.
\end{enumerate}

%Use $\perp$ and the We will proceed by proving that under an appropriate setting, such isomorphisms will be $\perpp$-isomorphisms, and so induce the $T$-homeomorphism $\phi:Y\to X$. However the same proof applies to $T^{-1}:\mathcal{A}(Y)\to\mathcal{A}(X)$, and the $T^{-1}$-homeomorphism will be $\phi^{-1}$.
%
%By proving that $T$ is basic we obtain the $T$-transform $\chi$. However the same proof will also apply to $T^{-1}$, which is therefore basic and this means precisely, by Proposition \ref{propositionbasicisomorphisms}(a) and (b), that the sections $\chi(y,\cdot)$ are injective and in the case that $\mathcal{A}(Y)$ is regular, item (d) implies that $\chi(y,\cdot)$ is also surjective. Proposition \ref{propositionmodelmorphism} will then allow us to classify isomorphisms (for a given signature) between $\mathcal{A}(X)$ and $\mathcal{A}(Y)$ in terms of isomorphisms between $H_X$ and $H_Y$. We will make no further mention of these facts.
%
%The general technique which we apply on each of these results can be thought as follows: First describe the relation $\perp$ in terms of the given algebraic structure on the space of continuous functions. Then, given $f$ and $g$, characterize functions $\widetilde{f}$ which ``resemble'' the characteristic function of $\supp(f)$, and then use $f\perpp g\iff\widetilde{f}\perp g$, however the way in which a function might ``resemble'' a characteristic function will depend on what algebraic structure we're considering.
%
%After that, given a classification of isomorphisms (for a given signature) between $H_X$ and $H_Y$, we prove that $T$ 

\subsection{Milgram's Theorem}

In this subsection, we will always assume that $X$ is a locally compact Hausdorff space, $H=\mathbb{R}$ and $\theta=0$ (the zero map from $X$ to $\mathbb{R}$). We denote $C_c(X,\mathbb{R})=C_c(X,0)$ (as in Definition \ref{definitionsigma}). If $Y$ is another locally compact space then $C_c(Y,\mathbb{R})$ is treated similarly. Note that $C_c(X,\mathbb{R})$ is regular by Urysohn's Lemma.

First, we present a well-known classification of continuous multiplicative isomorphisms of $\mathbb{R}$ (see \cite[Lemma 4.3]{MR0029476}, for example). We present its proof for the sake of completeness. Given $t\in\mathbb{R}$, $\operatorname{sgn}(t)$ denotes the sign of $t$ (the sign of $0$ is $\operatorname{sgn}(0)=0$).%, i.e.,
%\[\operatorname{sgn}(t)=\begin{cases}1&\text{, if }t>0,\\
%0&\text{, if }t=0,\\
%-1&\text{, if }t<0.\end{cases}\]

\begin{proposition}\label{theoremclassificationmultiplicativeisomorphismsofr}
Let $\tau:\mathbb{R}\to\mathbb{R}$ be a multiplicative isomorphism. Then
\begin{enumerate}
\item[(a)] Given $x\in\mathbb{R}$, $x\geq 0$ if and only if $\tau(x)\geq 0$;
\item[(b)] $\tau(-x)=-\tau(x)$ for all $x\in\mathbb{R}$;
\item[(c)] The following are equivalent:
\begin{enumerate}
\item[(1)] $\tau$ is continuous;
\item[(2)] $\tau$ is continuous at $0$;
\item[(3)] If $0<x<1$ then $0<\tau(x)<1$;
\item[(4)] $\tau$ is increasing;
\item[(5)] $\tau$ has the form $\tau(x)=\operatorname{sgn}(x)|x|^p$ for some $p>0$;
\end{enumerate}
\end{enumerate}
\end{proposition}
\begin{proof}
\begin{enumerate}
\item[(a)] Simply note that $x\geq 0$ if and only if $x=y^2$ for some $y$.
\item[(b)] If $x=0$ this is trivial. If $x\neq 0$, then $-x$ is the only number satisfying $(-x)^2=x^2$ and $-x\neq x$.
\item[(c)] The implications (5)$\Rightarrow$(1)$\Rightarrow$(2) are trivial.

\noindent(2)$\Rightarrow$(3): If $0<x<1$ then $x^n\to 0$, so $\tau(x)^n\to \tau(0)=0$ which implies $\tau(x)<1$.

\noindent(3)$\Rightarrow$(4) is immediate from (a) and (b).

\noindent(4)$\Rightarrow$(5): Letting $p=\log_2(\tau(2))>0$ (because $\tau(2)>\tau(1)=1$), we have $\tau(2^q)=(2^q)^p$ for all $q\in\mathbb{Q}$. Thus the restriction of $\tau$ to $[0,\infty)$ is an increasing map with dense image, hence surjective and continuous. Moreover, $\tau$ coincides with $x\mapsto x^p$ on a dense set of $[0,\infty)$, thus $\tau(x)=x^p$ for all $x\geq 0$ and thus for all $x$ by (b).\qedhere
\end{enumerate}
\end{proof}

We will now classify multiplicative isomorphisms from $C_c(X,\mathbb{R})$ and $C_c(Y,\mathbb{R})$. Following the procedure outlined in page \pageref{generalprocedureforclassification}, we first describe $\perpp$ in multiplicative terms.

\begin{lemma}\label{lemmamilgram}
If $f,g\in C_c(X,\mathbb{R})$, then
\[f\perpp g\iff \exists h\in C_c(X,\mathbb{R})\text{ such that }hf=f\text{ and }hg=0\ntag\label{conditionmilgram}\]
\end{lemma}
\begin{proof}
Indeed, first assume $f\perpp g$, i.e., $\supp(f)\cap\supp(g)=\varnothing$. By Urysohn's Lemma, there exists $h\in C_c(X,\mathbb{R})$ such that $h=1$ on $\supp(f)$ and $\supp(h)\subseteq X\setminus\supp(g)$. Then $h$ satisfies the condition (\ref{conditionmilgram}).

Conversely, if there is $h\in C_c(X,\mathbb{R})$ satisfying (\ref{conditionmilgram}), then $h=1$ on $\supp(f)$ and so $f\Subset h$. As $hg=0$ if and only if $h\perp g$, we conclude that $f\perpp g$.\qedhere
\end{proof}
%. Both $C_c(X)$ and $C_c(Y)$ are regular, whereas disjointness is described by the multiplicative structure as
%\[f\perp g\iff fg=0\iff \forall h\in C_c(X), (fg)h=(fg).\]
%Moreover, we have
%\[f\perpp g\iff \exists h(hf=f\text{ and }h\perp g)\tag{Mi}\]
%To see this, note that if $hf=f$ if and only if $h=1$ whenever $f\neq 0$. If $f\perpp g$, then such $h$ exists by Urysohn's Lemma. Conversely, if such $h$ exists then $f\Subset h$ and $h\perp g$, which implies $f\perpp g$.
%It is a known fact that every continuous multiplicative automorphism of the real numbers has the form $t\mapsto |t|^p$ for some $p>0$ (\cite{MR0029476}, Lemma 4.3). We give the proof here for the sake of completeness.

\begin{theorem}[{Milgram's Theorem, \cite[Theorem A]{MR0029476}, for locally compact spaces}]
Let $X$ and $Y$ be locally compact Hausdorff spaces and let $T:C_c(X,\mathbb{R})\to C_c(Y,\mathbb{R})$ be a multiplicative isomorphism. Then there exists a homeomorphism $\phi:Y\to X$, a subset $F\subseteq Y$ whose points are isolated in $Y$ and such that $F\cap K$ is finite for any compact $K\subseteq Y$, and a continuous positive function $p:Y\setminus F\to(0,\infty)$ satisfying
\[Tf(y)=\operatorname{sgn}(f(\phi(y)))|f(\phi (y))|^{p(y)}\]
for all $f\in C_c(X,\mathbb{R})$ and $y\in Y\setminus F$.
\end{theorem}
\begin{proof}
By Lemma \ref{lemmamilgram}, $T$ is a $\perpp$-isomorphism, so let $\phi$ be the $T$-homeomorphism. We prove that $T$ is $\phi$-basic in a few steps:
\begin{enumerate}
\item If $f,h\in C_c(X,\mathbb{R})$, then $f=1$ on $\supp(h)$ if and only if $Tf=1$ on $\supp(Th)$:

To verify this, simply note that $f=1$ on $\supp(h)$ if and only if $fh=h$, and similarly for $Tf$. Since $T$ is multiplicative we are done.

\item If $f,h\in C_c(X,\mathbb{R})$, then $f\neq 0$ on $\supp(h)$ if and only if $Tf\neq 0$ on $\supp(Th)$:

If $f\neq 0$ on $\supp(h)$, we can find $g\in C_c(X,\mathbb{R})$ such that $g=1/f$ on $\supp(h)$. Item 1.\ implies that $TfTg=1$ on $\supp(Th)$, and in particular $Tf\neq 0$ on $\supp(Th)$.

\item If $f\in C_c(X,\mathbb{R})$ and $y\in Y$, then $f(\phi(y))\neq 0$ if and only if $Tf(y)\neq 0$:

Assume $f(\phi(y))\neq 0$. Choose $h\in C_c(X)$ such that $\phi(y)\in\supp(h)\subseteq[f\neq 0]$. Then item 2.\ implies that $Tf\neq 0$ on $\supp(Th)$, which contains $y$.

\item If $f,g,h\in C_c(X,\mathbb{R})$, then $f$ and $g$ coincide and are nonzero on $\supp(h)$ if and only if $Tf$ and $Tg$ coincide and are nonzero on $\supp(h)$:

This is an immediate consequence of item 1.%It also follows from (1) and (2) that if two functions $f$ and $g$ are non-zero and coincide on some $\sigma(h)$ then $Tf$ and $Tg$ coincide on $\sigma(Th)$.
\item In particular, from 3., $f(\phi(y))=0$ if and only if $Tf(y)=0$.
\item If $f\in C_c(X,\mathbb{R})$ and $y\in Y$, then $f(\phi(y))=1$ if and only if $Tf(y)=1$:

Suppose this was not the case, say $f(\phi(y))=1$ but $Tf(y)\neq 1$, and let us deduce a contradiction. Take a neighbourhood $Y'$ of $y$ with compact closure. In particular from 1., $f$ is not constant on any neighbourhood of $\phi(y)$, so there is a sequence of distinct points $y_n\in Y'$ and $r>0$ such that
\begin{enumerate}
\item[(i)] $f(y_n)^n\to 1$;
\item[(ii)] $|Tf(y_n)-1|>r$ for all $n$.
\end{enumerate}
By Propositions \ref{propositiondisjointopensets} and \ref{propositionconstructfunction}(b), we can consider a subsequence of $\left\{y_n\right\}_n$, if necessary, and take a continuous function $g:X\to\mathbb{R}$ which coincides with $f^n$ on a neighbourhood of $\phi(y_n)$ for each $n$. Property (i) allows us to consider another tail of the sequence $\left\{y_n\right\}_n$, if necessary, and change $g$ by $\max(g,1/2)$, so we may assume that $g\neq 0$ everywhere. Let $u\in C_c(X,\mathbb{R})$ be a function with $u=1$ on $\phi(Y')$, so $ug\in C_c(X,\mathbb{R})$ and $ug=f^n$ on a neighbourhood of $\phi(y_n)$ for all $n$.

Let $z$ be a cluster point of the sequence $\left\{y_n\right\}_n$, so in particular $T(ug)(z)$ is a cluster point of the sequence $T(ug)(y_n)=Tf(y_n)^n$, where this equality follows from 4. Since $|Tf(y_n)-1|>r>0$, the only possibility is $T(ug)(z)=0$, and by 5.\ this means that $0=(ug)(\phi(z))=g(\phi(z))$ (because $z\in\overline{Y'}$), contradicting the fact that $g$ is nonzero.
\end{enumerate}

We conclude, from 5.\ and 6.\, that $T$ is $\phi$-basic. Let $\chi:Y\times\mathbb{R}\to\mathbb{R}$ be the $T$-transform.

Let $F=\left\{y\in Y:\chi(y,\cdot)\text{ is discontinuous}\right\}$. We will now prove that $F\cap K$ is finite for any compact subset $K\subseteq Y$. Otherwise there would be distinct $y_1,y_2,\ldots\in K$ such that $\chi(y_n,\cdot)$ is unbounded on every neighbourhood of $0$, so we can consider a strictly decreasing sequence $t_n\to 0$ such that $\chi(y_n,t_n)>n$. Going to a subsequence if necessary, we can construct, by Proposition \ref{propositionconstructfunction}, $f\in C_c(X,\mathbb{R})$ with $f(\phi(y_n))=t_n$, so $Tf(y_n)>n$ for all $n$, a contradiction to the boundedness of $Tf$.

Now we show that $F$ is open: If $z\in F$, then there is $t\in(0,1)$ with $\chi(z,t)>1$. Take $f\in C_c(X,\mathbb{R})$ such that $f=t$ on a neighbourhood $U$ of $\phi(z)$. In particular $Tf(z)>1$, so there is a neighbourhood $W$ of $z$ such that $Tf>1$ on $W$. Then for all $y\in\phi^{-1}(U)\cap W$, $\chi(y,t)=Tf(y)>1$, so $y\in F$.

Therefore, $F$ consists of isolated points, since $Y$ is locally compact. For $y\not\in F$, Proposition \ref{propositionmodelmorphism} and Proposition \ref{theoremclassificationmultiplicativeisomorphismsofr} imply that $\chi(y,\cdot)$ has the form $\chi(y,t)=\operatorname{sgn}(t)|t|^{p(y)}$ for some $p(y)>0$. Let $U$ be any open subset of $Y$ not intersecting $F$ and with compact closure. Take any function $f\in C_c(X,\mathbb{R})$ with $f=2$ on $\phi(U)$. Then for $y\in U$,
\[2^{p(y)}=\chi(y,2)=\chi(y,f(\phi(y)))=Tf(y)\]
so $p(y)=\log_2(Tf(y))$ for $y\in U$, showing that $p$ is continuous on $U$. Since $Y\setminus F$ is the union of such $U$ we are done.\qedhere
\end{proof}

\subsection{Group-Valued functions; Hernández-Ródenas Theorem}

Given topological groups $G$ and $H$, denote by $\operatorname{AbsIso}(G,H)$ the set of algebraic group isomorphisms from $G$ to $H$, and by $\operatorname{TopIso}(G,H)$ the set of topological (i.e., homeomorphisms which are) group isomorphisms.

Let $X$ and $Y$ be compact Hausdorff spaces and $G$ a Hausdorff topological group. In \cite[Theorem 3.7]{MR2324919}, Hernández and Ródenas classified non-vanishing group morphisms (not necessarily \emph{isomorphisms}) $T:C(X,G)\to C(Y,G)$ which satisfy the following properties:
\begin{enumerate}
\item[(i)] There exists a continuous group morphism $\psi:G\to C(X,G)$, where $C(X,G)$ is endowed with the topology of pointwise convergence, such that for all $\alpha\in G$ and all $y\in Y$, $T(\psi(\alpha))(y)=\alpha$;
\item[(ii)] For every continuous endomorphism $\theta:G\to G$ and every $f\in C(X,G)$, $T(\theta\circ f)=\theta\circ (Tf)$.
\end{enumerate}
If $T$ is a group isomorphism and $T^{-1}$ is continuous (with respect to uniform convergence) then condition (i) is immediately satisfied, however this is not true for (ii): For example, if $\operatorname{TopIso}(G,G)$ is non-abelian and $\rho\in\operatorname{TopIso}(G,G)$ is any non-central element, then the map $T:C(X,G)\to C(X,G)$ given by $Tf=\rho\circ f$ is a group isomorphism, and a self-homeomorphism of $C(X,G)$ with the topology of uniform convergence, which does not satisfy (ii).

In the next theorem we obtain the same type of classification as in \cite[Theorem 3.7]{MR2324919}, without assuming condition (ii), however we consider only non-vanishing group isomorphisms. Given a compact Hausdorff space $X$, a topological group $G$ and $\alpha\in G$, denote by $\overline{\alpha}$ the constant function $X\to G$, $x\mapsto \alpha$. We endow $C(X,G)$ with the topology of pointwise convergence.

\begin{theorem}%[{Compare with \cite[Theorem 3.7]{MR2324919}}]
Suppose that $G$ and $H$ are Hausdorff topological groups, and $X$ and $Y$ are compact Hausdorff spaces for which $(X,\overline{1_G},C(X,G))$ and $(Y,\overline{1_H},C(Y,H))$ are regular.

Let $T:C(X,G)\to C(Y,H)$ be a non-vanishing group isomorphism (Definition \ref{definitionnonvanishingbijection}). Then there exist a homeomorphism $\phi:Y\to X$ and a map $w:Y\to\operatorname{AbsIso}(H,G)$ such that $Tf(y)=w(y)(f(\phi(y))$ for all $y\in Y$ and $f\in C(X,G)$.

If $T$ is continuous on the constant functions then each $w(y)$ is continuous and $T$ is continuous on $C(X,G)$. If both $T$ and $T^{-1}$ are continuous on the constant functions then $w(y)\in\operatorname{TopIso}(H,G)$ and $T$ is a homeomorphism for the topologies of pointwise convergence.
\end{theorem}
\begin{proof}
By Theorem \ref{theoremnonvanishing} and Proposition \ref{propositionbasicnessofgroupvalued}, there is a homeomorphism $\phi:Y\to X$ such that $T$ is $\phi$-basic, and the sections $\chi(y,\cdot):G\to H$ of the $(\phi,T)$-transform $\chi$ are group morphisms by Proposition \ref{propositionmodelmorphism}. Similar facts hold for $T^{-1}$, so Proposition \ref{propositionbasicisomorphisms} implies that each section $\chi(y,\cdot)$ is bijective. Letting $w(y)=\chi(y,\cdot)$ we are done with the first part.

Now note that for all $\alpha\in G$ and $y\in Y$,
\[w(y)(\alpha)=w(y)(\overline{\alpha}(\phi(y)))=T\overline{\alpha}(y)\]
which implies that every $w(y)$ is continuous if and only if $T$ is continuous on the constant functions (because the map
\[G\to\left\{\overline{\alpha}:\alpha\in G\right\}\subseteq C(X,G),\qquad\alpha\mapsto\overline{\alpha}\]
is a homeomorphism). In this case, from the equality
\[Tf(y)=w(y)(f(\phi(y))\qquad\text{for all}\quad f\in C(X,G)\quad\text{and}\quad y\in Y,\]
we can readily see that $T$ is continuous. The last part, assuming also that $T^{-1}$ is continuous on the constant functions, is similar, using $T^{-1}$ and $w(y)^{-1}$ in place of $T$ and $w(y)$.\qedhere
\end{proof}

%%%%%%%%%%%%%%%%%%%%%%%%%%%%%%%%%%%%%%%%%%%%%%%%%%%%%%%%%%%%%%%%%%%%%%%%%%%%%%%%%%%%%%%%%%%%%%%%%%%%%%%%%%%%%%%%%%%%%%%%%%%%%%%%%%%%%%%%%%%%%%%%%%%%%%%%%%%%%%%%%%%%%%%%%%%%%%%%%%%%%%%%%%%%%%%%%%%%%%%%%%%%%%%%%%%%%%%%%%%%%%%%%%%%%%%%%%%%%%%%%%%%%%%%%%%%%%%%%%%%
\subsection{Kaplansky's Theorem}\label{subsectionkaplansky}

Let $R$ be a totally ordered set without supremum or infimum, considered as a topological space with the order topology, and let $X$ be a locally compact Hausdorff space. We consider the pointwise order on $C(X,R)$: $f\leq g$ if and only if $f(x)\leq g(x)$ for all $x$, which makes $C(X,R)$ a lattice: for all $f,g\in C(X,R)$ and $x\in X$,
\[(f\lor g)(x)=\max(f(x),g(x)),\qquad\text{and}\qquad(f\land g)(x)=\min(f(x),g(x)).\]
Denote by $C_b(X,R)$ the sublattice of bounded continuous functions from $X$ to $R$.

%\begin{proposition}
%Suppose $X$ is $R$-separated and $R$ does not contain a maximum nor a minimum. Then for every pair $F,G$ of disjoint compact subset of $X$ and $\alpha,\beta\in R$, there is $f\in C_b(X,R)$ which is equal to $\alpha$ on $F$ and $\beta$ on $G$.
%\end{proposition}
%\begin{proof}
%WLOG $\alpha<\beta$. Assume first that $G=\left\{x\right\}$ is a point. Using compactness of $F$ and $R$-separatedness, we can find functions $f_1,\ldots,f_n$ such that
%\begin{itemize}
%\item $f_i(x)=\beta$;
%\item $F\subseteq\bigcup_{i=1}^n f_i^{-1}(-\infty,\alpha)$;
%\end{itemize}
%and therefore $f=\max(\min(f_1,\ldots,f_n),\alpha)$ has the desired properties.
%
%The general case follows from the case above with a similar argument.\qedhere
%\end{proof}

We will consider sublattices $\mathcal{A}$ of $C_b(X,R)$ which satisfy
\begin{enumerate}
\item[(L1)] for all $f,g\in\mathcal{A}$, $\overline{[f\neq g]}$ is compact;
\item[(L2)] for all $f\in\mathcal{A}$, every open set $U\subseteq X$, every $x\in U$ and $\alpha\in R$, there exists $g\in\mathcal{A}$ such that $g(x)=\alpha$ and $[g\neq f]\subseteq U$.
\end{enumerate}

\begin{example}[{Kaplansky, \cite{MR0020715}}]\label{examplekaplanskyisl1l2}
Suppose that $X$ is compact and $\mathcal{A}$ is an $R$-normal\ftnt{Following \cite{MR0020715}, this means that for any pair of disjoint closed sets $F,G\subseteq X$ and $\alpha,\beta\in R$, there is $f\in\mathcal{A}$ which equals $\alpha$ on $F$ and $\beta$ on $G$.} sublattice of $C(X,R)$. Condition (L1) is trivial, so let us check that $\mathcal{A}$ satisfies (L2): Suppose $f$, $U$, $x$ and $\alpha$ are as in that condition. For the sake of the argument we can assume $f(x)\leq\alpha$. Let $\beta$ be any lower bound of $f(X)$, and from $R$-normality find $h\in\mathcal{A}$ such that $h(x)=\alpha$ and $h=\beta$ outside $U$. Then $g=f\lor h$ has the desired properties.
\end{example}

%In the next example, we consider the context of \cite{MR3162258}.

\begin{example}[{Li--Wong,  \cite{MR3162258}}]
Suppose that $X$ is compact, $R=\mathbb{R}$, and $\mathcal{A}$ is a regular additive subgroup of $C(X,\mathbb{R})$, so (L1) is also trivial and again we need to verify condition (L2): Let $f$, $U$, $x$ and $\alpha$ be as in (L2). By regularity, take $h$ such that $\supp(h)\subseteq U$ and $h(x)=\alpha-f(x)$. Then $g=f+h$ has the desired properties
\end{example}

We will now recover the main result of \cite{MR0020715}, The following lemma is based on \cite{MR0052760}.

\begin{lemma}\label{importantlemmakaplansky}
Let $\mathcal{A}$ be a sublattice of $C_b(X,R)$ satisfying (L1) and (L2), and let $f_0$ be any element of $\mathcal{A}$. Let $\mathcal{A}_{\geq f_0}=\left\{f\in\mathcal{A}:f\geq f_0\right\}$. Then $(X,f_0,\mathcal{A}_{\geq f_0})$ is weakly regular and for $f,g\in\mathcal{A}_{\geq f_0}$,
\begin{enumerate}
\item[(a)] $f\perp g\iff f\land g=f_0$ (which is the minimum of $\mathcal{A}_{\geq f_0})$;
\item[(b)] $f\Subset g$ is equivalent to the following statement:
\[\begin{split}
\text{``for every bounded subset $\H\subseteq\mathcal{A}$ such that $h\subseteq f$ for all $h\in \H$,}\\
\text{there is an upper bound $k$ of $\H$ such that $k\subseteq g$.''}\end{split}\tag{K}\]
\end{enumerate}
\end{lemma}

\begin{proof}
Weak regularity is immediate from (L2) and the fact that $R$ does not have a supremum, and item (a) is trivial. Let us prove (b). First suppose $f\Subset g$ and $\H\subseteq\mathcal{A}$ is a bounded subset such that $h\subseteq f$ for all $h\in \H$. Let $\alpha\in R$ be an upper bound of $\bigcup_{h\in\H}h(X)$. From weak regularity and compactness of $\supp(f)$, we can take finitely many functions $k_1,\ldots,k_n$ such that $k_i\Subset g$, and for every $x\in\supp(f)$ there is some $i$ with $k_i(x)>\alpha$. Letting $k=\bigvee_{i=1}^n k_i$ we obtain the desired properties.

Conversely, suppose that condition (K) holds. Let $\alpha$ be any upper bound of $f_0(X)$ and again take $\beta>\alpha$. Let $\mathscr{H}=\left\{h\in\mathcal{A}_{\geq f_0}:h\leq \beta, h\subseteq f\right\}$. Let $k$ be an upper bound of $\H$ with $k\subseteq g$. By Property (L2), we have $\sigma(f)=\bigcup_{h\in\mathscr{H}}\sigma(h)$, so $k\geq\beta$ on $\sigma(f)$ and thus also on $\supp(f)$, which implies $f\Subset k$. Since $k\subseteq g$ then $f\Subset g$.\qedhere
\end{proof}

As an immediate consequence of Lemma \ref{importantlemmakaplansky} and Theorem \ref{maintheorem} we have the following generalization of Kaplansky's Theorem:

\begin{theorem}[Kaplansky \cite{MR0020715}]\label{theoremkaplansky}
Suppose $R$ has no supremum nor infimum, $\mathcal{A}(X)$ and $\mathcal{A}(Y)$ are sublattices of $C_b(X,R)$ and $C_b(Y,R)$, respectively, satisfying conditions (L1) and (L2), and $T:\mathcal{A}(X)\to\mathcal{A}(Y)$ is a lattice isomorphism. Then for any $f_0\in\mathcal{A}$, $T$ restricts to a $\perpp$-isomorphism of the regular sublattices $\mathcal{A}(X)_{\geq f_0}$ and $\mathcal{A}(Y)_{\geq Tf_0}$. In particular, $X$ and $Y$ are homeomorphic.
\end{theorem}

Our immediate goal is to prove that the homeomorphism between $X$ and $Y$ given by Theorem \ref{theoremkaplansky} does not depend on the choice of function $f_0$ in Lemma \ref{importantlemmakaplansky}.

%\begin{example}
%In general, lattices isomorphisms between lattices of functions are not basic. Let $X=Y=\left\{x,y\right\}$ be a discrete space with two points, and $R=\mathbb{R}$ the real line with the usual order. Defining $T:C(X,\mathbb{R})\to C(Y,\mathbb{R})$ by
%\[Tf=\begin{cases}
%2f&\text{, if }0<f(x)\text{ and }0<f(y)\\
%f&\text{, otherwise},\end{cases}\]
%we obtain a lattice isomorphism. we obtain a non-basic lattice isomorphism. Indeed, if $f\leq g$, then there are two possibilities:
%\begin{cases}
%\item[Case 1:] $0<g(x)$ and $0<g(y)$;
%
%Then $Tf\leq\max(f,2f)\leq \max(g,2g)=2g=Tg$;
%
%\item[Case 2:] $g(x)\leq 0$ or $g(y)\leq 0$.
%
%Then $f(x)\leq 0$ or $f(y)\leq 0$ as well, so $Tf=f\leq g=Tg$.
%\end{cases}
%Therefore $T$ is a lattice isomorphism, and . Denoting by $\chi_z$ the characteristic function of $\left\{x\right\}$, we have $\chi_x(x)=(\chi_x+\chi_y)(x)$, but $T\chi_x(x)

\begin{lemma}\label{lemmarestrictionoflatticeisomorphismisbasic}\index{$T$-homeomorphism}
Under the conditions of Theorem \ref{theoremkaplansky}, let $\phi:Y\to X$ be the $T|_{\mathcal{A}(X)_{\geq f_0}}$-homeomorphism. If $f,g\in\mathcal{A}(X)_{\geq{f_0}}$ then $\phi^{-1}(\operatorname{int}([f=g]))=\operatorname{int}([Tf=Tg])$.
\end{lemma}
\begin{proof}
We will use the superscript ``$f_0$'' as in Definition \ref{definitionsigma}, that is, for all $f\geq f_0$.
\[\sigma^{f_0}(f)=\operatorname{int}(\overline{[f\neq f_0]}),\qquad\text{and}\qquad Z^{f_0}(f)=\operatorname{int}([f=f_0]).\]
and similarly with $Tf_0$ in place of $f_0$. Then $\phi$ is the only homeomorphism satisfying $\sigma^{f_0}(f)=\phi(\sigma^{Tf_0}(Tf))$, or equivalently $Z^{f_0}(f)=\phi(Z^{Tf_0}(Tf))$, for all $f\geq f_0$.

First, assume that $f\leq g$, and let $U=\operatorname{int}([f=g])$. For all $x\in[f<g]$, choose a function $k_x\in\mathcal{A}(X)_{\geq f_0}$ such that
\begin{itemize}
\item $\sigma^{f_0}(k_x)\cap [f=g]=\varnothing$;
\item $k_x(x)=g(x)$;
\item $k_x\leq g$.
\end{itemize}
Then $g=\sup\left\{f,k_x:x\in[f<g]\right\}$, and since $T$ is a lattice isomorphism we obtain $Tg=\sup\left\{Tf,Tk_x:x\in[f<g]\right\}$.

Let us prove that $Tf(x)=Tg(x)$ for all $x\in\phi^{-1}(U)$.
%
Since $U\subseteq\bigcap_{x\in[f<g]}Z^{f_0}(k_x)$, then $\phi^{-1}(U)\subseteq\bigcap_{x\in[x<g]}Z^{Tf_0}(Tk_x)$.

Given $y\in\phi^{-1}(U)$, use property (L2) to find $h\in\mathcal{A}(Y)$ such that $h(y)=Tf(y)$ and $h=Tg$ outside of $\phi^{-1}(U)$. Then $h'=(Tf\lor h)\land Tg$ is an upper bound of $\left\{Tf,Tk_x:x\in[f<g]\right\}$, so it is also an upper bound of $Tg$, and in particular,
\begin{align*}
Tg(y)\leq ((h\lor Tf)\land Tg)(y)=Tf(y)\land Tg(y)\leq Tf(y)
\end{align*}
so $Tg(y)=Tf(y)$.

In the general case, if $f=g$ on an open set $U$ then $f=f\land g=g$ on $U$, so the previous case implies that $Tf$ and $Tg$ coincide with $T(f\land g)$ on $\phi^{-1}(U)$.\qedhere
\end{proof}

\begin{theorem}\label{theoremkaplanskywithoutlowerbound}
Under the conditions of Theorem \ref{theoremkaplansky}, there exists a unique homeomorphism $\phi:Y\to X$ such that $\phi(\operatorname{int}([Tf=Tg]))=\operatorname{int}([f=g])$ for all $f,g\in\mathcal{A}(X)$. (In this case we will still call $\phi$ the \emph{$T$-homeomorphism}.)
\end{theorem}
\begin{proof}
For all $f_0\in \mathcal{A}(X)$, let $\phi^{f_0}:Y\to X$ be the $T|_{\mathcal{A}(X)_{\geq f_0}}$-homeomorphism. Given $f_0,g_0\in\mathcal{A}(X)$, Lemma \ref{lemmarestrictionoflatticeisomorphismisbasic} implies that $\phi^{f_0\land g_0}$ satisfies the property of the $T|_{\mathcal{A}(X)_{\geq f_0}}$-homeomorphism, which uniquely defines $\phi^{f_0}$, and similarly for $\phi^{g_0}$, so $\phi^{f_0}=\phi^{f_0\land g_0}=\phi^{g_0}$. We are done by letting $\phi=\phi^{f_0}$ for some arbitrary $f_0\in\mathcal{A}(X)$.\qedhere
\end{proof}

A natural goal now is to classify the lattice isomorphisms as given in Theorem \ref{theoremkaplanskywithoutlowerbound}. Although at the moment we cannot prove that all lattice isomorphisms as above are basic, this is true in some cases. A similar argument to the one used in the proof of Theorem \ref{propositionbasicisomorphisms} appears in \cite{MR2995073}, although in a different context (considering lattices of possibly unbounded real-valued continuous functions on complete metric spaces). See also \cite{MR2999998,MR3404615}.

\begin{theorem}\label{theoremclassificationlatticemorphismfirstcountable}
Suppose that $X$ and $Y$ are first-countable (and locally compact Hausdorff), that $R=\mathbb{R}$ with the usual order, and that $T:C_c(X,0)\to C_c(Y,0)$ is a lattice isomorphism. Then there are a unique homeomorphism $\phi:Y\to X$ and a continuous function $\chi:Y\times\mathbb{R}\to\mathbb{R}$ such that
\[Tf(y)=\chi(y,f(\phi(y)))\qquad\text{ for all }y\in Y\text{ and }f\in C_c(X,0)\ntag\label{equationclassificationoforderkaplanskyisomorphismsforfirstcountble}\]
and $\chi(y,\cdot):\mathbb{R}\to\mathbb{R}$ is an increasing bijection for all $y\in Y$.
\end{theorem}
\begin{proof}
Suppose $x\in X$ and $f(x)=g(x)$. Using Propositions \ref{propositiondisjointopensets} and \ref{propositionconstructfunction}(b), we can find a sequence $\left\{x_n\right\}_n$ of points in $X$ converging to $x$ and a function $h\in C_c(X,0)$ such that, for all $n$, $h=f$ on some neighbourhood of $x_{2n}$ and $h=g$ on some neighbourhood of $x_{2n+1}$. Therefore, $Tf$ will coincide with $Th$ on some neighbourhood of $\phi^{-1}(x_{2n})$ for all $n\in\mathbb{N}$,
\[Th(\phi^{-1}(x))=\lim_{n\to\infty}Th(\phi^{-1}(x_{2n}))=\lim_{n\to\infty}Tf(\phi^{-1}(x_{2n}))= Tf(\phi^{-1}(x))\]
and similarly $Tg(\phi^{-1}(x))=Th(\phi^{-1}(x))=Tf(\phi^{-1}(x))$, which proves that $T$ is basic.

Therefore, there is a function $\chi:Y\times \mathbb{R}\to\mathbb{R}$ such that Equation (\ref{equationclassificationoforderkaplanskyisomorphismsforfirstcountble}) is satisfied. Every section $\chi(y,\cdot):\mathbb{R}\to\mathbb{R}$ is surjective, since $C_c(Y,0)$ is regular, and the same argument with $T^{-1}$, and Proposition \ref{propositionbasicisomorphisms} imply that $\chi(y,\cdot)$ is also injective, hence a bijection.

Proposition \ref{propositionmodelmorphism}, applied to the signature of lattices (with the binary symbol ``$\lor$'' interpreted as ``join'') implies that the sections $\chi(y,\cdot)$ are lattice isomorphisms of $\mathbb{R}$ for all $n$, and in particular homeomorphisms. Theorem \ref{theorembasicisomorphisms} implies that $\chi$ is continuous (and $T$ is a homeomorphism for the topologies of pointwise convergence.)
\end{proof}

\begin{example}
There are non-first countable spaces for which the conclusion of Theorem \ref{propositionbasicisomorphisms} holds.

Let $\Omega=\omega_1\cup\left\{\omega_1\right\}$ be the successor of the first uncountable ordinal. We extend\ftnt{The fact that $\omega_1\not\in\omega_1$ follows from the Axiom of Regularity, see \cite[p.\ 3]{MR1940513}.} the order of $\omega_1$ to $\Omega$ by setting $\alpha<\omega_1$ for all $\alpha\in\omega_1$, and $\Omega$ is a compact Hausdorff space with the order topology. If $Z$ is any first-countable space, then any continuous function $f:\Omega\to Z$ is constant on a neighbourhood of $\omega_1$: Indeed, for every $n\in\mathbb{N}$ choose $\alpha_n<\omega_1$ such that $|f(\beta)-f(\omega_1)|<1/n$ whenever $\beta\geq\alpha_n$. Letting $\alpha=\sup_n\alpha_n$, we have $\alpha<\omega_1$ and $f(\beta)=f(\omega_1)$ for all $\beta\in[\alpha,\omega_1]$

Now suppose that $R=\mathbb{R}$ and $T:C(\Omega)\to C(\Omega)$ is a lattice isomorphism, and let $\phi:\Omega\to\Omega$ be the $T$-homeomorphism. Since $\omega_1$ is the only non-$G_\delta$ point of $\Omega$, we have $\phi(\omega_1)=\omega_1$. The previous paragraph allows us to identify the lattices
\[C_c(\omega_1)\simeq\left\{f\in C(\Omega):f(\omega_1)=0\right\},\]
which then induces a lattice isomorphism $T|_{\omega_1}:C_c(\omega_1)\to C_c(\omega_1)$. In this case, note that $\phi|_{\omega_1}$ is the $T|_{\omega_1}$-homeomorphism. We can now prove that $T$ is basic with respect to $\phi$.

Let $f,g\in C(\Omega)$. If $f(\omega_1)=g(\omega_1)$, then the first paragraph implies that $f=g$ on some neighbourhood of $\omega_1$ and thus $Tf(\omega_1)=Tg(\omega_1)$. If $f(\alpha)=g(\alpha)$ for some $\alpha<\omega_1$, consider $\widetilde{f}\in C_c(\omega_1,0)$ given by
\[\widetilde{f}(x)=\begin{cases}
f(x),&\text{if }x\leq\alpha,\\
0,&\text{otherwise,}\end{cases}\]
and define $\widetilde{g}$ similarly. Since $(-\infty,\alpha]$ is open in $\Omega$, then
\[Tf(\phi^{-1}(\alpha))=T\widetilde{f}(\phi^{-1}(\alpha))\qquad\text{and}\qquad Tg(\phi^{-1}(\alpha))=T\widetilde{g}(\phi^{-1}(\alpha)),\]
and since $\omega_1$ is first-countable, we use the previous example to conclude that
\[T\widetilde{f}(\phi^{-1}(\alpha))=T\widetilde{g}(\phi^{-1}(\alpha))\]
so $Tf(\phi^{-1}(\alpha))=Tg(\phi^{-1}(\alpha))$. Therefore $T$ is basic (with respect to $\phi$).
\end{example}

%\begin{theorem}\label{theoremgeneralkaplansky}
%Let $R=\mathbb{R}$ and $T:C_c(X,0)\to C_c(Y,0)$ be a lattice isomorphism. Then there exists a unique homeomorphism $\phi:Y\to X$ and a unique continuous map $\xi:Y\times R\to R$ such that each section $\xi(y,\cdot):R\to R$ is an order automorphism of $R$ (and in particular it is a homeomorphism) and $Tf(y)=\chi(y,f(\phi(y))$ for all $y\in Y$ and $f\in C_c(X,0)$.
%\end{theorem}
%\begin{proof}
%Let $\phi$ be the $T$-homeomorphism. We prove that $T$ is basic with respect to $\phi$. Suppose $f,g\in C_c(X,0)$, $x\in X$ and $f(x)=g(x)$. If $x$ is isolated then $f=g$ on the open set $\left\{x\right\}$ so $Tf=Tg$ on $\left\{\phi^{-1}(x)\right\}$.

%If $x$ is not isolated, we can find a sequence of distinct points $\left\{x_n\right\}_n$ such that $Tf(x_n)\to Tf(x)$ and $Tconsider a converging net $y_i\to x$ with $y_i\neq x$ for all $x$. By transfinite induction, we can find a subnet $\alpha\mapsto y{i(\alphaj)}$ (indexed by some limit ordinal) such that $y_{i(\alpha)}\neq y_{i(\beta)}$ whenever $\alpha\neq\beta$. Since $(f(x_{i(\alpha)}),g(x_{i(\alpha)}))\to(f(x),g(x))$, we can find a 

%Given $f_0\in\mathcal{A}(X)$, Lemma \ref{lemmarestrictionoflatticeisomorphismisbasic} implies that $T|_{\mathcal{A}(X)_{\geq f_0}}$ is a basic $\perpp$-isomorphism, so there are a homeomorphism $\phi^{f_0}:Y\to X$ and a partial map $\chi^{f_0}:\dom(\chi^{f_0})\subseteq Y\times R\to R$ for which
%\[Tf(y)=\chi^{f_0}(y,f(\phi^{f_0}(y)))\]
%whenever $f\geq f_0$. Moreover, lattices have a signature with a binary function symbol $\lor$ which is interpreted in the usual manner: $a\lor b$ is simply the supremum of $a$ and $b$. By Proposition \ref{propositionmodelmorphism}, the sections $\chi^{f_0}(y,\cdot):R\to R$ ($y\in Y$) are lattice morphisms of $R$.
%
%Now we prove that $\phi^{f_0}$ and $\chi^{f_0}$ are actually independent of $f_0$.
%
%Suppose $f_0\leq g_0$ in $\mathcal{A}(X)$. By Lemma \ref{lemmarestrictionoflatticeisomorphismisbasic}, we have, for all $g\geq g_0$,
%\[\phi^{f_0}\left(\operatorname{int}\left(\overline{[Tf\neq Tg_0]}\right)\right)=\operatorname{int}\left(\overline{[f\neq g_0]}\right)\]
%which is the characterizing property of $\phi^{g_0}$. Therefore $\phi^{f_0}=\phi^{g_0}$.
%
%Given $y\in Y$ and $\alpha\in R$, choose $g\geq g_0$ such that $g(\phi(y))=\alpha$. Then
%\[\chi^{f_0}(y,\alpha)=\chi^{f_0}(y,g(\phi(y)))=Tg(y)=\chi^{g_0}(y,g(\phi(y)))=\chi^{g_0}(y,\alpha)\]
%and this proves that $\chi^{f_0}=\chi^{g_0}$.
%
%Now for the general case, when $f_0$ and $g_0$ are non-comparable, we simply use the function $\min(f_0,g_0)$ and the previous case to conclude that
%\[\phi^{f_0}=\phi^{\min(f_0,g_0)}=\phi^{g_0}\qquad\text{and}\qquad\chi^{f_0}=\chi^{\min(f_0,g_0)}=\chi^{g_0}.\qedhere\]
%\end{proof}

%%%Let $R=\mathbb{R}$ and $\theta=0$. By a \emph{regular additive sublattice} of $C_c(X)$ we mean a collection $\mathcal{A}\subseteq C_c(X)$ which is both regular, an additive subgroup, and a sublattice of $C_c(X)$.

In the case additive lattice isomorphisms of spaces of real-valued functions, we do not require the first-countability hypothesis.

\begin{theorem}\label{theoremadditivekaplansky}
Suppose $R=\mathbb{R}$, and $T:C_c(X)\to C_c(Y)$ is an additive lattice isomorphism. Then there are a unique homeomorphism $\phi:Y\to X$ and a unique positive continuous function $p:Y\to (0,\infty)$ such that $Tf(y)=p(y)f(\phi(y))$ for all $f\in C_c(X)$ and $y\in Y$.
\end{theorem}
\begin{proof}
%%%By Theorem \ref{theoremgeneralkaplansky}, there are a function $\chi:Y\times\mathbb{R}\to\mathbb{R}$ with increasing sections $\chi(y,\cdot)$ ($y\in\mathbb{R}$) and a homeomorphism $\phi:Y\to X$ such that $Tf(y)=\chi(y,f(\phi(y)))$ for all $y\in Y$ and $f\in\mathcal{A}(X)$. Proposition \ref{propositionmodelmorphism} implies that each section $\chi:(y,\cdot):\mathbb{R}\to\mathbb{R}$ is additive.

First note that for all $f\in C_c(X)$, $|f|=(f\lor 0)-(f\land 0)$, so $T|f|=|Tf|$. Let $\phi:Y\to X$ be the $T$-homeomorphism, given by Theorem \ref{theoremkaplanskywithoutlowerbound}.

Now suppose $f(x)=0$ but $Tf(\phi^{-1}(x))\neq 0$. First take a compact neighbourhood $U$ of $x$ and $r>0$ such that $|Tf|>r$ on $\phi^{-1}(U)$. Moreover, $f$ is not constant on any neighbourhood of $x$, so there is a sequence of distinct points $x_n\in U$ such that $|f(x_n)|<n^{-2}$. 
 Using Propositions \ref{propositiondisjointopensets} and \ref{propositionconstructfunction}(b), we can take a subsequence if necessary and consider $g\in C_c(X)$ such that for all $n$, $g=nf$ on a neighbourhood of $x_n$. Then $Tg=nTf$ on a neighbourhood of $\phi^{-1}(x_n)$, however
 \[nr<nTf(\phi^{-1}x_n)=nTg(\phi^{-1}x_n),\]
 which contradicts the fact that $g$ is bounded.
%By Proposition \ref{propositiondisjointopensets}, take a subsequence if necessary and disjoint neighbourhoods $U_n$ of each $x_n$ such that $|f|<n^{-2}$ on $U_n$. Then for each $n$ take, by regularity, $g_n\in\mathcal{A}$ such that $g_n(x_n)> n|f_n(x_n)|$ and $\supp(g_n)\subseteq U_n$. Let $h_n=g_n\land nf_n$.

%Then in fact the family $\left\{h_n:n\in\mathbb{N}\right\}$ is bounded ($h_n\leq 1$ for all $n$) and $\supp(h_n)\subseteq\overline{U}$, and it is easy enough, using regularity and compactness of $U$ to construct $k\in\mathcal{A}$ with $k\geq 1$ on $\overline{U}$, so $\left\{Th_n:n\in\mathbb{N}\right\}$ is also bounded by $Tk$. Then for all $n$, $h_n=nf_n$ on a neighbourhood of $x_n$, so
%\[nr<nTf(\phi^{-1}(x_n))=Th(\phi^{-1}(x_n))\leq k(\phi^{-1}(x_n))\leq\sup_{y\in Y}k(y)\]
%contradicting the fact that $k$ is bounded.

Therefore $T$ is basic with respect to $\phi$, so let $\chi$ be the $T$-transform. Each section $\chi(y,\cdot)$ is an additive order-preserving bijection (Propositions \ref{propositionmodelmorphism} and \ref{propositionbasicisomorphisms}) and hence has the form $\chi(y,t)=p(y)t$ for some $p(y)>0$. If $Tf(y)\neq 0$, then $f(\phi(y))\neq 0$ as well and $p=Tf/(f\circ \phi)$ on a neighbourhood of $y$, thus $p$ is continuous.\qedhere
%: The signature of lattices contains a binary function symbol $\lor$, namely $x\lor y=\max(x,y)$, and increasing maps are precisely $\lor$-morphisms. Since $T$ is a lattice morphism, each section $\chi(y,\cdot)$ is also a lattice morphism and in particular preserves $\lor$, or equivalently is increasing.
%
%Thus, for a given $y\in Y$, $\chi(y,\cdot)$ is an increasing, additive bijection, %, hence it is continuous
%so it has the form $\chi(y,t)=p(y)t$ for a certain (unique) $p(y)=\chi(y,1)>0$. If $Tf(y)\neq 0$, then $f(\phi y)\neq 0$ as well and $p=Tf(/f\circ \phi)$ on a neighbourhood of $y$, so $p$ is continuous.\qedhere
\end{proof}

%%%%%%%%%%%%%%%%%%%%%%%%%%%%%%%%%%%%%%%%%%%%%%%%%%%%%%%%%%%%%%%%%%%%%%%%%%%%%%%%%%%%%%%%%%%%%%%%%%%%%%%%%%%%%%%%%%%%%%%%%%%%%%%%%%%%%%%%%%%%%%%%%%%%%%%%%%%%%%%%%%%%%%%%%%%%%%%%%%%%%%%%%%%%%%%%%%%%%%%%%%%%%%%%%%%%%%%%%%%%%%%%%%%%%%%%%%%%%%%%%%%%%%%%%%%%%%%%%%%%

\subsection{Li--Wong Theorem}\label{subsectionliwong}

%, i.e., supportsThroughout this subsection, we let $X$ and $Y$ be compact Hausdorff spaces, $C=\mathbb{R}$, $\theta=0$ and $\mathcal{A}(X)$ and $\mathcal{A}(Y)$ vector sublattices of $C(X)$ and $C(Y)$, respectively, which are regular (or \emph{completely regular} in the nomenclature of \cite{MR3162258}).% We also assume that $\mathcal{A}(X)$ and $\mathcal{A}(Y)$ contain all constant functions.

In \cite{MR3162258}, Li and Wong proved Theorem \ref{theoremliwong}, which can be seen as a generalization of Theorem \ref{theoremadditivekaplansky}. We will proceed in the opposite direction, i.e., by proving their result (or more precisely, the specific case where the domains are compact) instead as a consequence of the more general Theorem \ref{theoremkaplansky}.

\begin{theorem}[Li--Wong \cite{MR3162258}]\label{theoremliwong}
Let $X$ and $Y$ be compact Hausdorff spaces, and $\mathcal{A}(X)$ and $\mathcal{A}(Y)$ be two regular vector sublattices of $C(X,\mathbb{R})$ and $C(Y,\mathbb{R})$, respectively. Suppose that $T:\mathcal{A}(X)\to \mathcal{A}(Y)$ is a linear isomorphism which preserves non-vanishing functions, that is, for all $f\in\mathcal{A}(X)$,
\[0\in f(X)\iff 0\in Tf(Y).\]
Then there is a homeomorphism $\phi:Y\to X$ and a continuous non-vanishing function $p:Y\to\mathbb{R}$ such that $Tf(y)=p(y)f(\phi(y))$ for all $f\in\mathcal{A}(X)$ and $y\in Y$.
\end{theorem}
\begin{proof}
In order to apply Theorem \ref{theoremkaplansky}, we need to modify $T$ to obtain a lattice isomorphism. Since $\mathcal{A}(X)$ is a sublattice, then for all $f\in\mathcal{A}(X)$,
\[f^+=\max(f,0),\quad f^-=\max(-f,0)\quad\text{and}\quad|f|=f^++f^-\text{ belong to }\mathcal{A}(X).\]

As $\mathcal{A}(X)$ is regular, we can take finitely many functions $f_1,\ldots,f_n\in\mathcal{A}(X)$ such that for all $x\in X$, $f_i(x)\neq 0$ for some $i$, and therefore $F=\sum_{i=1}^n|f_i|\in\mathcal{A}(X)$ and $F$ is non-vanishing, so $TF$ is also non-vanishing. We define new classes of functions
\[\mathcal{B}(X)=\left\{f/F:f\in\mathcal{A}(X)\right\},\qquad\text{and}\qquad\mathcal{B}(Y)=\left\{f/TF:f\in\mathcal{A}(Y)\right\}\]
It is immediate to see that $\mathcal{B}(X)$ and $\mathcal{B}(Y)$ are regular, and that $\mathcal{B}(X)$ is a sublattice of $C(X,\mathbb{R})$ because $F>0$. However we need to prove that $\mathcal{B}(Y)$ is a sublattice of $C(Y,\mathbb{R})$.

Let $U=[TF>0]$ and $V=[TF<0]$, which are complementary clopen sets in $Y$, and $1_U$ and $1_V$ their respective characteristic functions. We first prove that for all $f\in\mathcal{A}(Y)$, we have $f1_U\in\mathcal{A}(Y)$. Let $k=\sup_{y\in Y}|f(y)|$.

Use regularity of $\mathcal{A}(Y)$, compactness of $U$ and the fact that $\mathcal{A}(Y)$ is a sublattice to find a function $u\in\mathcal{A}(X)$ such that $u>k$ on $U$ and $u=0$ on $V$, and similarly $v\in\mathcal{A}(Y)$ with $v<-k$ on $V$ and $v=0$ on $U$. Then
\[f1_U=(f\land (u-v))\lor 0\in\mathcal{A}(Y)\]
It also follows that $f1_V=f-f1_U\in\mathcal{A}(Y)$. Therefore if $b\in\mathcal{B}(Y)$ then $b(TF)$, $b(TF)1_U$ and $b(TF)1_V$ belong to $\mathcal{A}(Y)$, thus
\[(b\lor 0)=\frac{(b(TF)1_U\lor 0)}{TF}+\frac{(b(TF)1_V\land 0)}{TF}\in\mathcal{B}(Y).\]
In general, if $b,c\in\mathcal{B}(Y)$ then
\[b\lor c=b+((c-b)\lor 0)\in\mathcal{B}(Y),\]
and therefore $\mathcal{B}(Y)$ is a sublattice of $C(Y,\mathbb{R})$.

Now define a linear isomorphism $S:\mathcal{B}(X)\to\mathcal{B}(Y)$, $S(f)=T(fF)/TF$, which preserves non-vanishing functions. Note that in this case, the constant functions at $1$ on $X$ and $Y$ belong to $\mathcal{B}(X)$ and $\mathcal{B}(Y)$, and that $S(1)=1$. Given a scalar $\lambda$, linearity and the non-vanishing property of $S$ imply that, for all $f\in\mathcal{B}(X)$,
\begin{align*}
\lambda\not\in f(X)&\iff f-\lambda\text{ is non-vanishing}\\
&\iff Sf-\lambda\text{ is non-vanishing}\iff \lambda\not\in Sf(Y),
\end{align*}
so $f(X)=Sf(Y)$ and in particular $f\geq 0$ if and only if $Sf\geq 0$. Therefore $S$ is a lattice isomorphism, and we can apply Kaplansky's theorem (\ref{theoremkaplanskywithoutlowerbound}) to find the $S$-homeomorphism $\phi:Y\to X$. Now we need to prove that $S$ is $\phi$-basic.

Suppose $f(x)\neq 0$ for a given $x\in X$, and let us assume, without loss of generality, that $f(x)>0$. Then $f>0$ on some neighbourhood $U$ of $x$. Again using compactness of $X\setminus U$ and regularity of the sublattice $\mathcal{B}(X)$ we can construct a function $g\in\mathcal{B}(X)$ such that $g=0$ on some neighbourhood of $x$ and $g>0$ on $X\setminus U$. Letting $\widetilde{f}=f\lor g$, we have $\widetilde{f}=f$ on some neighbourhood of $x$, so $S\widetilde{f}=Sf$ on some neighbourhood of $\phi^{-1}(x)$. But $\widetilde{f}$ is non-vanishing, so $S\widetilde{f}$ is also non vanishing and in particular $Sf(\phi^{-1}(x))\neq 0$. This proves that $S$ is basic with respect to $\phi$. Letting $\chi:Y\times\mathbb{R}\to\mathbb{R}$ be the $S$-transform, we have that all sections $\chi(y,\cdot)$ are linear and increasing (Theorem \ref{propositionmodelmorphism}), hence of the form
\[\chi(y,t)=P(y)t\]
for a certain $P(y)>0$. Denoting the constant function $x\mapsto 1$ by $1$ (either on $X$ or $Y$), then
\[P(y)=\chi(y,1)=\chi(y,1(\phi(y)))=S1(y)=1(y)=1.\]
that is, $\chi(y,t)=t$ for all $t\in\mathbb{R}$.

Finally, for all $f\in\mathcal{A}(X)$ and $y\in Y,$
\[Tf(y)=(TF)(y)\left[S\left(\frac{f}{F}\right)(y)\right]=(TF)(y)\chi\left(y,\frac{f}{F}(\phi(y))\right)=\frac{TF(y)}{F(\phi(y))}f(\phi(y)),\]
as we wanted.\qedhere
%we see that Since $T(1)\neq 0$ everywhere, dividing $T$ by $T(1)$ if necessary (and changing $\mathcal{A}(Y)$ by $\mathcal{A}(Y)/T(1)$) we may assume $T(1)=1$. By linearity and since $T$ preserves non-vanishing functions, $f(X)=Tf(Y)$, i.e., $T$ preserves images. In particular $f\geq 0$ if and only if $Tf\geq 0$, which actually allows us to assume the real case. Using \ref{theoremadditivekaplansky} and multiplying again by $T(1)$ we obtain $Tf(y)=p(y)f(\phi y)$ for suitable $p$ and $\phi$, and it is immediate to see that $p=T(1)$ in this case.\qedhere
\end{proof}

We can extend this result to the complex case as follows:

\begin{corollary}\label{corollaryliwongcomplex}
Suppose $T:C(X,\mathbb{C})\to C(Y,\mathbb{C})$ is a linear isomorphism which preserves non-vanishing functions. Then there is a homeomorphism $\phi:Y\to X$ and a continuous non-vanishing function $p:Y\to\mathbb{C}$ such that $Tf(y)=p(y)f(\phi(y))$ for all $f\in C(X,\mathbb{C})$ and $y\in Y$.
\end{corollary}
\begin{proof}
Substituting $T$ by $S=T/T(1)$ if necessary, we may assume $T(1)=1$. The same argument as in the proof above implies $Tf(Y)=f(X)$ for all $f\in C(X,\mathbb{C})$, and in particular the restriction of $T$ to $C(X,\mathbb{R})$ gives us a linear isomorphism $T|_{C(X,\mathbb{R})}:C(X,\mathbb{R})\to C(Y,\mathbb{R})$ preserving non-vanishing functions. We apply Theorem \ref{theoremliwong} to this case to obtain $p$ and $\phi$ satisfying $Tf=p\cdot (f\circ\phi)$ for all $f\in C(X,\mathbb{R})$. Since $C(X,\mathbb{R})$ generates $C(X,\mathbb{C})$ then the same formula is valid for all $f\in C(X,\mathbb{C})$.\qedhere
\end{proof}
%%%%%%%%%%%%%%%%%%%%%%%%%%%%%%%%%%%%%%%%%%%%%%%%%%%%%%%%%%%%%%%%%%%%%%%%%%%%%%%%%%%%%%%%%%%%%%%%%%%%%%%%%%%%%%%%%%%%%%%%%%%%%%%%%%%%%%%%%%%%%%%%%%%%%%%%%%%%%%%%%%%%%%%%%%%%%%%%%%%%%%%%%%%%%%%%%%%%%%%%%%%%%%%%%%%%%%%%%%%%%%%%%%%%%%%%%%%%%%%%%%%%%%%%%%%%%%%%%%%%

\subsection{Jarosz' Theorem}\label{subsectionjarosz}

Throughout this subsection, we fix $\mathbb{K}=\mathbb{R}$ or $\mathbb{C}$. Given a locally compact Hausdorff space $X$, we let $C_c(X)$ be the Banach space of $\mathbb{K}$-valued, compactly supported, continuous function on $X$, where supports are the usual ones, i.e., $\supp(f)=\overline{[f\neq 0]}$.

%\begin{theorem}\label{lemmaadditivelattice}
%Let $T:C_c(X)\to C_c(Y)$ be an additive $\perp$-isomorphism, and let $h \in C_c(X)$ be arbitrary. Then $f(x)=0$ for some point $x\in\supph$ if and only if $Tf(y)=0$ for some point $y\in\suppTh$.

%In particular, if $T$ is a $\perpp$-isomorphism then $T$ is basic.
%\end{theorem}
%\begin{proof}
%Suppose $Tf(x)\neq 0$ for all $x\in\suppTh$. In particular, $f$ is not constant $0$ on any open subset of $\sigma(h)$, and so there is a sequence of distinct points $x_n\in\sigma(h)$ such that $nf(x_n)\to 0$, and for some $r>0$ we have $|Tf|>r$ on $\sigma(Th)$. By \ref{lemmadisjointopensets} and \ref{lemmaconstructfunction} we can, up to taking a subsequence, construct a function $g\in C_c(X)$ such that $g=nf$ on a neighbourhood of $x_n$, so $|Tg(y_n)|\geq nr$ for all $n$, a contradiction to $Tg$ being bounded.

%If $T$ is a $\perpp$-isomorphism, let $\phi$ be the $T$-homeomorphism and suppose $f(\phi y)=0$. Then the closure of every neighbourhood of $y$ has some point $z$ such that $Tf(z)=0$. Since $Y$ is locally compact Hausdorff (and $Tf$ is continuous) this implies $Tf(y)=0$.\qedhere
%\end{proof}

\begin{theorem}[Jarosz \cite{MR1060366}]\label{theoremjarosz}
If $T:C_c(X)\to C_c(Y)$ is a linear $\perp$-isomorphism, then there exist a homeomorphism $\phi:Y\to X$ and a continuous non-vanishing function $p:Y\to\mathbb{K}$ such that $Tf(y)=p(y)f(\phi(y))$ for all $f\in C_c(X)$ and $y\in Y$.
\end{theorem}
\begin{proof}
\textbf{First assume that $X$ and $Y$ are compact}, and let us show that $f\neq 0$ everywhere if and only if $Tf\neq 0$ everywhere. Suppose otherwise, say $f(x)=0$, and we have two cases: first, if $f$ is constant on a neighbourhood of $x$, this means that $Z(f)\neq \varnothing$, and Theorem \ref{theoremdisjoint} implies that $Z(Tf)\neq \varnothing$, and in particular $Tf(y)=0$ for any $y\in Z(Tf)$.

In the second case, if $f$ is not constant on any neighbourhood of $x$, we proceed in the same manner as in Theorem \ref{theoremadditivekaplansky}: first find a sequence of distinct point $x_n\in X$ such that $0<|f(x_n)|<n^{-2}$. Propositions \ref{propositiondisjointopensets} and \ref{propositionconstructfunction}(b) allows us to take a subsequence if necessary and construct $g\in C(X)$ which coincides with $nf$ on a neighbourhood of $x_n$ for all $n$. However, linearity and the $\perp$-preserving property imply that $Tg=nTf$ on some open set as well (Theorem \ref{theoremdisjoint}), which contradicts $Tf$ being bounded because $\sup_{y\in Y}|Tf(y)|>0$.%  and $\overline{U_n}\cap\overline{\bigcup_{m\neq n}U_m}=\varnothing$ for all $n$. Define $g$ as $nf$ on $U_n$ and extend it arbitrarily to a continuous function on $X$. We have $g=nf$ on an some open set, so by linearity and since $T$ preserves disjointness, $Tg=nTf$ on some open set as well (Theorem \ref{theoremdisjoint}). However $\inf_{y\in Y}|Tf(y)|>0$, contradicting $Tg$ being bounded.

The result follows in this case from the Li--Wong Theorem (Theorem \ref{theoremliwong} or Corollary \ref{corollaryliwongcomplex}).

\textbf{Now let $X$ and $Y$ be arbitrary locally compact Hausdorff.} Given $b\in C_c(X)$, set $T_b:C(\supp(b))\to C(\supp(Tb))$ as $T_bf=(Tf')|_{\supp(Tb)}$, where $f'$ is any element of $C_c(X)$ extending $f$. Note that $T_bf$ does not depend on the choice of $f'$, since, for all $f',g'\in C_c(X)$,
\begin{align*}
f'|_{\supp(b)}=g'|_{\supp(b)}&\iff\sigma(b)\subseteq [f'=g']\iff\sigma(b)\subseteq Z(f'-g')\\
&\iff\sigma(b)\cap\sigma(f'-g')=\varnothing\iff b\perp(f'-g'),
\end{align*}
and the last condition is preserved by $T$ since it is an additive $\perp$-isomorphism (Theorem \ref{theoremdisjoint}).

Since $f\perpp g$ if and only if $f|_{\supp(b)}\perpp g|_{\supp(b)}$ for all $b$, the previous case allows us to obtain functions $p^b$ and $\phi^b$ such that $Tf(y)=p^b(y)f(\phi^b(y))$ for all $y\in\supp(Tb)$. Clearly, if $b\subseteq b'$ then $p^{b'}|_{\supp(b)}=p^b$ and $\phi^{b'}|_{\supp(b)}=\phi^{b'}$. Thus defining $p$ and $\phi$ as $p(y)=p^b(y)$ and $\phi(y)=\phi^b(y)$ where $b\in C_c(X)$ is such that $y\in\supp(b)$ we obtain the desired maps.\qedhere

%In particular $T(1)\neq 0$ everywhere. Dividing $T$ by $T(1)$ if necessary, we may assume first $T(1)=1$. Then by the previous paragraph, $T$ preserves images, i.e., $f(X)=Tf(Y)$ for all $f\in C(X)$, and in particular $T$ preserves positive functions and the order on those. Therefore the same condition (K) (as in Kaplansky's theorem) describes strong disjointness.

%Therefore $T$ is a $\perpp$-isomorphism, so take a homeomorphism $\phi$ as in \ref{maintheorem}. The proof that $T$ follows from the first paragraph: If $f(x)=0$ but $Tf(\phi^{-1}(x))\neq 0$, then it is easy to take a function $g$ such that $Tg$ coincides with $Tf$ on a neighbourhood of $\phi^{-1}(x)$, and $Tg\neq 0$ everywhere. However $g$ will then coincide with $f$ on a neighbourhood of $x$, so $g(x)=0$, a contradiction with the first part.

%Thus $T$ is basic. The same argument as in the proof of Banach-Stone gives us the desired form for $T$.\qedhere
\end{proof}

%%%%%%%%%%%%%%%%%%%%%%%%%%%%%%%%%%%%%%%%%%%%%%%%%%%%%%%%%%%%%%%%%%%%%%%%%%%%%%%%%%%%%%%%%%%%%%%%%%%%%%%%%%%%%%%%%%%%%%%%%%%%%%%%%%%%%%%%%%%%%%%%%%%%%%%%%%%%%%%%%%%%%%%%%%%%%%%%%%%%%%%%%%%%%%%%%%%%%%%%%%%%%%%%%%%%%%%%%%%%%%%%%%%%%%%%%%%%%%%%%%%%%%%%%%%%%%%%%%%%

\subsection{Banach-Stone Theorem}\label{subsectionbanachstone}

We use the same notation as in the previous subsection. Given a locally compact Hausdorff space $X$, endow $C_c(X)$ with the supremum norm: $\Vert f\Vert_\infty=\sup_{x\in X}|f(x)|$.%We fix $\mathbb{K}=\mathbb{R}$ or $\mathbb{K}=\mathbb{C}$. Given a locally compact Hausdorff space $X$, we let $C_c(X)$ be the Banach of $\mathbb{K}$-valued, compactly supported continuous function on $X$, endowed with the supremum norm: $\Vert f\Vert_\infty=\sup_{x\in X}|f(x)|$.

Recall that, by the Riesz-Markov-Kakutani Representation Theorem (\cite[Theorem 2.14]{MR924157}), continuous linear functionals on $C_c(X)$ correspond to (integration with respect to) regular Borel measures on $X$. As a consequence, the extremal points $T$ of the unit ball of the dual of $C_c(X)$ have the form $T(f)=\lambda f(x)$ for some $x\in X$ and $|\lambda|=1$.

%Suppose $\mu$ is an extremal measure with $\Vert\mu\Vert=1$, so that $|\mu|$ is actually a probability measure, and $\mu=\int gd|\mu|$ for some $g$. We show that $g$ is constant a.e. If this was not the case, there would be a partition $\mathbb{C}=A\cup B$ such that $g^{-1}(A)$ and $g^{-1}(B)$ have nonzero measure. Let $g_C$ be the function which is equal to $g$ on $g^{-1}(C)$ and zero everywhere else. Then $g=g_A+g_B$, and
%\[\mu=\int g_Ad|\mu|+\int g_Bd|\mu|=\Vert g_A\Vert_{1,|\mu|} \mu_A+\Vert g_B_{1,|\mu|}\Vert\mu_B\]
%where $\mu_C=(\int g_Cd|\mu|)\Vert g_C\Vert_{1,|\mu|}$ and $\Vert\cdot\Vert_{1,|\mu|}$ denotes the $L^1$ norm with respect to $|\mu|$. Since $g_A$ and $g_B$ have disjoint supports, we get a contradiction to extremality of $\mu$.

%Therefore $g$ is constant, and it is easy to see that $|g|=\Vert\mu\Vert=1$, so we can assume $\mu=|\mu|$. Now we show that $\mu$ has the desired form. The same type of argument shows that for every subset $A\subseteq X$, $\mu(A)=0$ or $1$. If every point had measure $0$, then regularity of $\mu$ would give us an open cover of null open sets, so every compact set would have measure zero as well and $\mu$ would be zero on $C_c(X)$ and therefore on $C_0(X)$ by continuity, a contradiction. Thus there is some point $x$ with $\mu(x)=1$.

Given $f\in C_c(X)$, denote by $N(f)$ the set of extremal points $T$ in the unit ball of the dual space $C_c(X)^*$ such that $T(f)\neq 0$. From the previous paragraph we obtain
\[f\perp g\iff N(f)\cap N(g)=\varnothing\tag{BS},\]
and the Banach-Stone Theorem is an immediate consequence of Jarosz' Theorem.

\begin{theorem}[Banach-Stone \cite{MR1501905}]\label{theorembanachstone}
Let $X$ and $Y$ be locally compact Hausdorff spaces and let $T:C_c(X)\to C_c(Y)$ be an isometric linear isomorphism. Then there exists a homeomorphism $\phi:Y\to X$ and a continuous function $p:Y\to\mathbb{S}^1$ for which
\[Tf(y)=p(y)f(\phi(y))\qquad\forall f\in C(X),\ \forall y\in Y.\]
\end{theorem}
%\begin{proof}
%$f\perp g$ if and only if $N(f)\cap N(g)=\varnothing$, so the theorem follows from Jarosz's Theorem.\qedhere

%By (BS) $T$ is a $\perpp$-isomorphism, and \ref{maintheorem} gives us appropriate homeomorphism $\phi$ and function $\chi$.

%To obtain basicality of $T$, we need to show that $f(x)=0$ implies $f(\phi(x))=0$. Given $\epsilon>0$, define $n_\epsilon:\mathbb{C}\to\mathbb{C}$ by
%\[n_\epsilon(z)=\begin{cases}
%0&\text{, if }|z|\leq\epsilon\\
%\frac{z}{|z|}(|z|-\epsilon)&\text{, if }|z|\geq \epsilon.
%\end{cases}\]
%and $f_\epsilon=n_\epsilon\circ f$. Then $x\in\rho(f_\epsilon)$, so $\phi^{-1}(x)\in\rho(Tf_\epsilon)$ and in particular $Tf_\epsilon(\phi^ {-1}(x))=0$. Since $f_\epsilon\to f$ uniformly, then $Tf_\epsilon\to Tf$ uniformly as well so $Tf(\phi^{-1}(x))=0$.

%By \ref{theorembasicisomorphisms}, $T$ is of the form
%\[Tf(y)=\chi(y,f(\phi^{-1}(y)))\]

%for a suitable function $\chi$. Linearity of $T$ implies that $\chi(y,\cdot):\mathbb{K}\to\mathbb{K}$ is a linear function, hence of the form $\chi(y,t)=p(y)t$ for all $t\in\mathbb{K}$ for some number $p(y)\in\mathbb{K}$. Given $y$, we can choose a function $f\in C(X)$ for which $Tf(y)\neq 0$, and so, in a neighbourhood of $y$, $p$ is of the form
%\[p(y)=\frac{Tf(y)}{f(\phi^{-1}(y))}\]
%so $p:Y\to\mathbb{K}$ is continuous. That $|p(y)|=1$ follows from $T$ being isometric.\qedhere
%\end{proof}

%%%%%%%%%%%%%%%%%%%%%%%%%%%%%%%%%%%%%%%%%%%%%%%%%%%%%%%%%%%%%%%%%%%%%%%%%%%%%%%%%%%%%%%%%%%%%%%%%%%%%%%%%%%%%%%%%%%%%%%%%%%%%%%%%%%%%%%%%%%%%%%%%%%%%%%%%%%%%%%%%%%%%%%%%%%%%%%%%%%%%%%%%%%%%%%%%%%%%%%%%%%%%%%%%%%%%%%%%%%%%%%%%%%%%%%%%%%%%%%%%%%%%%%%%%%%%%%%%%%%
